% LPO_en_Haskell.tex
% Lógica de primer orden (LPO) en Haskell.
% Eduardo Paluzo
% Sevilla, 16 de julio de 2016
% =============================================================================

\documentclass[a4paper,12pt,twoside]{book}

%%%%%%%%%%%%%%%%%%%%%%%%%%%%%%%%%%%%%%%%%%%%%%%%%%%%%%%%%%%%%%%%%%%%%%%%%%%%%%
%%  Paquetes adicionales                                                   %%
%%%%%%%%%%%%%%%%%%%%%%%%%%%%%%%%%%%%%%%%%%%%%%%%%%%%%%%%%%%%%%%%%%%%%%%%%%%%%%

\usepackage[utf8x]{inputenc}       % Acentos de UTF8
\usepackage[spanish]{babel}        % Castellanización.
\usepackage[T1]{fontenc}           % Codificación T1 con European Computer
                                   % Modern.  
\usepackage{graphicx}
\usepackage{fancyvrb}              % Verbatim extendido
\usepackage{makeidx}               % Índice
\usepackage{amsmath}               % AMS LaTeX
\usepackage{amsthm} 
\usepackage{latexsym}

% Fuentes
\usepackage{mathpazo}              % Fuentes semejante a palatino
\usepackage[scaled=.90]{helvet}
\usepackage{cmtt}
\renewcommand{\ttdefault}{cmtt}
\usepackage{a4wide}
% \usepackage{xmpincl}               % Incluye metadato de licencia CC.

% Tikz
\usepackage{tkz-berge}

% Desactivar <,> cuando se hacen dibujos con tikz.
\tikzset{
every picture/.append style={
  execute at begin picture={\deactivatequoting},
  execute at end picture={\activatequoting}
  }
}

% Márgenes
\usepackage[margin=1in]{geometry}

% Control de espacios en la tabla de contenidos:
\usepackage[titles]{tocloft}
\setlength{\cftbeforechapskip}{2ex}
\setlength{\cftbeforesecskip}{0.5ex}
\setlength{\cftsecnumwidth}{12mm}
\setlength{\cftsubsecindent}{18mm}

% Control de listas
% Elimina espacio entre item y párrafo y coloca item en el margen izquierdo
% \usepackage{enumitem}
% \setlist[enumerate,itemize]{noitemsep, nolistsep, leftmargin=*}

\usepackage{minitoc}

% Doble espacio entre líneas
\usepackage{setspace}
\onehalfspacing

% \linespread{1.05}                % Distancia entre líneas
\setlength{\parindent}{5mm}        % Indentación de comienzo de párrafo
\deactivatetilden                  % Elima uso de ~ para la eñe
\raggedbottom                      % No ajusta los espacios verticales.

\usepackage[%
  pdftex,
  pdfauthor={Eduardo Paluzo},%
  pdftitle={LPO en Haskell},%
  pdfstartview=FitH,%
  bookmarks=false,%
  colorlinks=true,%
  urlcolor=blue,%
  unicode=true]{hyperref}      

\usepackage{tikz}

%%%%%%%%%%%%%%%%%%%%%%%%%%%%%%%%%%%%%%%%%%%%%%%%%%%%%%%%%%%%%%%%%%%%%%%%%%%%%%
%%  Cabeceras                                                              %%
%%%%%%%%%%%%%%%%%%%%%%%%%%%%%%%%%%%%%%%%%%%%%%%%%%%%%%%%%%%%%%%%%%%%%%%%%%%%%%

\usepackage{fancyhdr}

\addtolength{\headheight}{\baselineskip}

\pagestyle{fancy}

\cfoot{}

\fancyhead{}
\fancyhead[RE]{\bfseries \nouppercase{\leftmark}}
\fancyhead[LO]{\bfseries \nouppercase{\rightmark}}
\fancyhead[LE,RO]{\bfseries \thepage}

%%%%%%%%%%%%%%%%%%%%%%%%%%%%%%%%%%%%%%%%%%%%%%%%%%%%%%%%%%%%%%%%%%%%%%%%%%%%%%
%%  Definiciones                                                           %%
%%%%%%%%%%%%%%%%%%%%%%%%%%%%%%%%%%%%%%%%%%%%%%%%%%%%%%%%%%%%%%%%%%%%%%%%%%%%%%

\input definiciones
\def\mtctitle{Contenido}

%%%%%%%%%%%%%%%%%%%%%%%%%%%%%%%%%%%%%%%%%%%%%%%%%%%%%%%%%%%%%%%%%%%%%%%%%%%%%%
%%  Título                                                                 %%
%%%%%%%%%%%%%%%%%%%%%%%%%%%%%%%%%%%%%%%%%%%%%%%%%%%%%%%%%%%%%%%%%%%%%%%%%%%%%%

\title{\Huge Lógica de primer orden en Haskell}
\author{Eduardo Paluzo}
\date{\vfill \hrule \vspace*{2mm}
  \begin{tabular}{l}
      \href{http://www.cs.us.es/glc}
           {Grupo de Lógica Computacional} \\
      \href{http://www.cs.us.es}
           {Dpto. de Ciencias de la Computación e Inteligencia Artificial} \\
      \href{http://www.us.es}
           {Universidad de Sevilla}  \\
      Sevilla, 16 de junio de 2016 (Versión de \today)
  \end{tabular}\hfill\mbox{}}

%%%%%%%%%%%%%%%%%%%%%%%%%%%%%%%%%%%%%%%%%%%%%%%%%%%%%%%%%%%%%%%%%%%%%%%%%%%%%%%
%%  Construcción del índice                                                 %%
%%%%%%%%%%%%%%%%%%%%%%%%%%%%%%%%%%%%%%%%%%%%%%%%%%%%%%%%%%%%%%%%%%%%%%%%%%%%%%%

\makeindex

%%%%%%%%%%%%%%%%%%%%%%%%%%%%%%%%%%%%%%%%%%%%%%%%%%%%%%%%%%%%%%%%%%%%%%%%%%%%%%
%%  Documento                                                              %%
%%%%%%%%%%%%%%%%%%%%%%%%%%%%%%%%%%%%%%%%%%%%%%%%%%%%%%%%%%%%%%%%%%%%%%%%%%%%%%

% \includeonly{Introduccion}

% \includexmp{licencia}

\begin{document}
\dominitoc

\maketitle
\newpage

\noindent
Esta obra está bajo una licencia Reconocimiento--NoComercial--CompartirIgual
2.5 Spain de Creative Commons. 

\vspace*{4ex}

\begin{center}
\fbox{
\begin{minipage}{175mm}
\vspace*{1ex}
\noindent                     
{\bf Se permite:}
\begin{itemize}                         
\item copiar, distribuir y comunicar públicamente la obra
\item hacer obras derivadas
\end{itemize}
         
\noindent       
{\bf Bajo las condiciones siguientes:}

\begin{minipage}[t]{20mm}
  \vspace*{0mm}\hspace*{1em}
  \includegraphics[scale=0.35]{deed.png}
\end{minipage}%
\begin{minipage}[t]{15cm}
  \vspace*{0mm}
  {\bf Reconocimiento}. Debe reconocer los créditos de la  obra de la manera
  especificada por el autor.  
\end{minipage}
\\[1ex]

\begin{minipage}[t]{20mm}
  \vspace*{0mm}\hspace*{1em}
  \includegraphics[scale=0.35]{deed-eu.png}
\end{minipage}%
\begin{minipage}[t]{15cm}
  \vspace*{2mm}
  {\bf No comercial}. No puede utilizar esta obra para fines comerciales.
\end{minipage}
\\[1ex]

\begin{minipage}[t]{20mm}
  \vspace*{0mm}\hspace*{1em}
  \includegraphics[scale=0.35]{deed_002.png}
\end{minipage}%
\begin{minipage}[t]{15cm}
  \vspace*{-2mm}
  {\bf Compartir bajo la misma licencia}. Si altera o transforma esta obra, o
  genera una obra derivada, sólo puede distribuir la obra generada bajo una
  licencia idéntica a ésta. 
\end{minipage}

\begin{itemize}    
\item Al reutilizar o distribuir la obra, tiene que dejar bien claro los
  términos de la licencia de esta obra. 
\item Alguna de estas condiciones puede no aplicarse si se obtiene el permiso
 del titular de los derechos de autor.
\end{itemize}
\vspace*{1ex}
\end{minipage}}
\end{center}

\vspace*{4ex}

Esto es un resumen del texto legal (la licencia completa). Para ver una copia
de esta licencia, visite 
\href{http://creativecommons.org/licenses/by-nc-sa/2.5/es/}
     {\url{http://creativecommons.org/licenses/by-nc-sa/2.5/es/}}
o envie una carta a Creative Commons, 559 Nathan Abbott Way, Stanford,
California 94305, USA.

%%% Local Variables: 
%%% mode: latex
%%% TeX-master: t
%%% End: 

\newpage 

\tableofcontents
\newpage

\chapter*{Introducción}
\addcontentsline{toc}{chapter}{Introducción}

El objetivo del trabajo es la implementación de los algoritmos de la lógica de
primer orden en Haskell. Consta de dos partes:

\begin{itemize}
\item La primera parte consiste en la implementación en Haskell de la teoría
  impartida en la asignatura
  \href{https://www.cs.us.es/~jalonso/cursos/lmf-15} 
       {Lógica matemática y fundamentos}\
  \footnote{\url{{https://www.cs.us.es/~jalonso/cursos/lmf-15}}}
  (\cite{Alonso-15a}). Para ello, se lleva a cabo la adaptación de los programas del libro de
  J. van Eijck 
  \href{http://citeseerx.ist.psu.edu/viewdoc/download?doi=10.1.1.467.1441&rep=rep1&type=pdf}
       {Computational semantics and type theory}\
  \footnote{\url{http://citeseerx.ist.psu.edu/viewdoc/download?doi=10.1.1.467.1441&rep=rep1&type=pdf}}
  (\cite{Eijck-03}) y su correspondiente teoría. 

\item En la segunda parte se comentan aquellos sistemas empleados en la elaboración del
  trabajo.
 
\end{itemize}



%%% Local Variables: 
%%% mode: latex
%%% TeX-master: "LPO_en_Haskell"
%%% End: 


\part{Semántica computacional y teoría de tipos}

\chapter{Programación funcional con Haskell}

En este capítulo se hace una breve introducción a la programación funcional en
Haskell suficiente para entender su aplicación en los siguientes
capítulos. Para una introducción más amplia se pueden consultar los apuntes de
la asignatura de Informática de 1º del Grado en Matemáticas
(\cite{Alonso-15b}). También se puede emplear como lectura complementaria, y se
ha empleado para algunas definiciones del trabajo (\cite{YerPal-90})

El contenido de este capítulo se encuentra en el módulo \texttt{PFH} 
\entrada{PFH}.

%%% Local Variables:
%%% mode: latex
%%% TeX-master: "LPO_en_Haskell"
%%% End:

\chapter{Sintaxis y semántica de la lógica de primer orden}

El \textbf{lenguaje de la lógica de primer orden} está compuesto por

\begin{enumerate}
\item Variables proposicionales $p,q,r,\dots$
\item Conectivas lógicas:
  \begin{center}
   \begin{tabular}{| l | l |}
     \hline
      $\neg$            & \text{Negación} \\ \hline
      $\vee$            & \text{Disyunción} \\ \hline
      $\wedge$          & \text{Conjunción} \\ \hline
      $\rightarrow$     & \text{Condicional} \\ \hline
      $\leftrightarrow$ & \text{Bicondicional}\\
     \hline
   \end{tabular}
 \end{center}
\item Símbolos auxiliares $"(",")"$
\item Cuantificadores: $\forall$ (Universal) y $\exists$ (Existencial)
\item Símbolo de igualdad: $=$
\item Constantes: $a,b,\dots,a_1,a_2,\dots$
\item Símbolos de relación: $P,Q,R,\dots $
\item Símbolos de función: $f,g,h\dots $
\end{enumerate}

\entrada{Modelo}.

\chapter{Deducción natural}

En este capítulo se pretende implementar la deducción natural de la lógica de
primer orden en Haskell. El contenido de este capítulo se encuentra en el
módulo \texttt{DNH}.  

\entrada{DNH}

\chapter{Prueba de teoremas en lógica de
        predicados}

Este capítulo pretende aplicar métodos de tableros para la demostración
de teoremas en lógica de predicados. El contenido de este capítulo
se encuentra en el módulo \texttt{PTLP}.
\entrada{PTLP}
\chapter{Modelos de Herbrand}

En este capítulo se pretende formalizar los modelos de Herbrand. Herbrand
propuso un método constructivo para generar interpretaciones de fórmulas o
conjuntos de fórmulas.

\entrada{HerbrandAlt}

\chapter{Resolución en lógica de primer orden}

En este capítulo se introducirá la resolución en la lógica de primer orden
y se implementará en Haskell. El contenido de este capítulo
se encuentra en el módulo \texttt{RES}

\entrada{RES}


\part{Sistemas utilizados}
\chapter{Apéndice: GitHub}
En este apéndice se pretende introducir al lector en
el empleo de \href{https://github.com/}{GitHub}, sistema remoto de versiones.
\section{Crear una cuenta}

El primer paso es crear una cuenta en la página web de \href{https://github.com/}{GitHub},
para ello \texttt{sign up} y se rellena el formulario.

\section{Crear un repositorio}

Mediante \texttt{New repository} se crea un repositorio nuevo.
Un repositorio es una carpeta de trabajo. En ella se guardarán todas
las versiones y modificaciones de nuestros archivos.

Necesitamos darle un nombre adecuado y seleccionar
\begin{enumerate}
\item En \texttt{Add.gitignore} se selecciona \texttt{Haskell}.
\item En \texttt{add a license} se selecciona \texttt{GNU General Public License v3.0}.
\end{enumerate}

Finalmente \texttt{Create repository}

\section{Conexión}

Nuestro interés está en poder trabajar de manera local y poder tanto actualizar
GitHub como nuestra carpeta local. Los pasos a seguir son
\begin{enumerate}
\item Generar una clave SSH mediante el comando
\begin{verbatim}
  ssh-keygen -t rsa -b 4096 -C ``tuCorreo''
\end{verbatim}
  Indicando una contraseña. Si queremos podemos dar una localización de guardado
  de la clave pública.
\item Añadir la clave a \texttt{ssh-agent}, mediante
\begin{verbatim}
  eval "$(ssh-agent -s)"
  ssh-add ~/.ssh/id_rsa
\end{verbatim}

\item Añadir la clave a nuestra cuenta. Para ello: \texttt{Setting} $\rightarrow$
  \texttt{SSH and GPG keys} $\rightarrow$ \texttt{New SSH key}. En esta última
  sección se copia el contenido de
\begin{verbatim}
~/.ssh/id_rsa.pub
\end{verbatim}
, por defecto. Podríamos haber puesto otra localización en el primer paso.
\item Se puede comprobar la conexión mediante el comando
\begin{verbatim}
ssh -T git@github.com
\end{verbatim}
\item Se introducen tu nombre y correo
\begin{verbatim}
 git config --global user.name "Nombre"
 git config --global user.email "<tuCorreo>"
\end{verbatim}
\end{enumerate}

\section{Pull y push}

Una vez hecho los pasos anteriores, ya estamos conectados con GitHub y podemos
actualizar nuestros ficheros. Nos resta aprender a actualizarlos.

\begin{enumerate}
\item Clonar el repositorio a una carpeta local:
  Para ello se ejecuta en una terminal
  \texttt{git clone <enlace que obtienes en el repositorio>}. El enlace que sale en el
  repositorio pinchando en \texttt{(clone or download)} y, eligiendo \texttt{(use SSH)}.

\item Actualizar tus ficheros con la versión de GitHub:
  Abierto emacs, se ejecuta \texttt{(Esc-x)+(magit-status)}.Para ver una lista de los cambios que
  están \texttt{(unpulled)}, se ejecuta en magit \texttt{remote update} . Se emplea \texttt{pull},
  y se actualiza. \texttt{(Pull: origin/master)}
\item Actualizar GitHub con tus archivos locales:
  Abierto emacs, se ejecuta \texttt{(Esc-x)+(magit-status)}. Sale la lista de los cambios
  \texttt{(UnStages)}. Con la \texttt{(s)} se guarda, la \texttt{(c)+(c)} hace \texttt{(commit)},
  Le ponemos un título y, finalmente \texttt{(Tab+P)+(P)} para hacer \texttt{(push)} y subirlos a GitHub.
\end{enumerate}

%%%%%%%%%%%%%%%%%%%%%%%%%%%%%%%%%%%%%%%%%%%%%%%%%%%%%%%%%%%%%%%%%%%%%%%%%%%%%%%
%%  Bibliografía                                                            %%
%%%%%%%%%%%%%%%%%%%%%%%%%%%%%%%%%%%%%%%%%%%%%%%%%%%%%%%%%%%%%%%%%%%%%%%%%%%%%%%

\addcontentsline{toc}{chapter}{Bibliografía}
\bibliographystyle{abbrv}
\bibliography{LPO_en_Haskell}

%%%%%%%%%%%%%%%%%%%%%%%%%%%%%%%%%%%%%%%%%%%%%%%%%%%%%%%%%%%%%%%%%%%%%%%%%%%%%%%
%%  Índice                                                                  %%
%%%%%%%%%%%%%%%%%%%%%%%%%%%%%%%%%%%%%%%%%%%%%%%%%%%%%%%%%%%%%%%%%%%%%%%%%%%%%%%

\addcontentsline{toc}{chapter}{Indice de definiciones}

\printindex

\end{document}



%%% Local Variables:
%%% mode: latex
%%% TeX-master: t
%%% End:
